% GENERATED FILE. Edit canonical lessons, then run npm run generate:books.
% canonical-source-hash: fnv1a64:67726e3e379f11c6
% canonical-lessons: SA-C103-kandukah, SA-C103-kridanakam, SA-C103-anguliyakam, SA-C103-kankanam, SA-C103-uttariyam

\chapter{Ball, Toy, Ring, Bangle, Shawl}
\label{ch:sa-ball-toy-ring-bangle-shawl}
\hlchaptermodality{103}

\begin{chapteropening}
\textbf{By the end of this chapter:} \emph{I can say a ball, a toy, a ring, a bangle and a shawl.}

Complete the last lesson of chapter 103: I can say a ball, a toy, a ring, a bangle and a shawl.
\end{chapteropening}

\section[\texorpdfstring{\sk{कन्दुकः} (kandukaḥ)}{कन्दुकः (kandukaḥ)}]{\sk{कन्दुकः} (kandukaḥ) — a ball}

\label{lesson:SA-C103-kandukah}



\begin{quote}
Before the new one: say the Sanskrit for a stove, then the Sanskrit for a candle.
\end{quote}

\subsection*{You'll want to know: \sk{कन्दुकः}}
\textbf{\sk{कन्दुकः}} — \emph{kandukaḥ} — \textquotedblleft{}a ball\textquotedblright{}.

\begin{grammarlens}[title={What you've built}]
One more word for this chapter.
\end{grammarlens}

\subsection*{Your turn}
\begin{itemize}
\raggedright
  \item \emph{Say it:} \emph{kandukaḥ}
  \item \emph{Say it:} \emph{kandukaḥ}, once more
  \item \emph{Say it:} say \emph{kandukaḥ}
  \item \emph{Recall:} say the Sanskrit for a stove, then the Sanskrit for a candle, then say \emph{kandukaḥ} again
\end{itemize}

\begin{tcolorbox}[breakable,colback=teal!4,colframe=teal!35!black,title={Before you move on}]
What does \sk{कन्दुकः} mean? (\textquotedblleft{}A ball\textquotedblright{}.) Say it once more.
\end{tcolorbox}
\section[\texorpdfstring{\sk{क्रीडनकम्} (krīḍanakam)}{क्रीडनकम् (krīḍanakam)}]{\sk{क्रीडनकम्} (krīḍanakam) — a toy}

\label{lesson:SA-C103-kridanakam}



\begin{quote}
Before the new one: say the Sanskrit for a candle, then the Sanskrit for a ball.
\end{quote}

\subsection*{You'll want to know: \sk{क्रीडनकम्}}
\textbf{\sk{क्रीडनकम्}} — \emph{krīḍanakam} — \textquotedblleft{}a toy\textquotedblright{}.

\begin{grammarlens}[title={What you've built}]
One more word for this chapter.
\end{grammarlens}

\subsection*{Your turn}
\begin{itemize}
\raggedright
  \item \emph{Say it:} \emph{krīḍanakam}
  \item \emph{Say it:} \emph{krīḍanakam}, once more
  \item \emph{Say it:} say \emph{krīḍanakam}
  \item \emph{Recall:} say the Sanskrit for a candle, then the Sanskrit for a ball, then say \emph{krīḍanakam} again
\end{itemize}

\begin{tcolorbox}[breakable,colback=teal!4,colframe=teal!35!black,title={Before you move on}]
What does \sk{क्रीडनकम्} mean? (\textquotedblleft{}A toy\textquotedblright{}.) Say it once more.
\end{tcolorbox}
\section[\texorpdfstring{\sk{अंगुलीयकम्} (aṅgulīyakam)}{अंगुलीयकम् (aṅgulīyakam)}]{\sk{अंगुलीयकम्} (aṅgulīyakam) — a ring}

\label{lesson:SA-C103-anguliyakam}



\begin{quote}
Before the new one: say the Sanskrit for a ball, then the Sanskrit for a toy.
\end{quote}

\subsection*{You'll want to know: \sk{अंगुलीयकम्}}
\textbf{\sk{अंगुलीयकम्}} — \emph{aṅgulīyakam} — \textquotedblleft{}a ring\textquotedblright{}.

\begin{grammarlens}[title={What you've built}]
One more word for this chapter.
\end{grammarlens}

\subsection*{Your turn}
\begin{itemize}
\raggedright
  \item \emph{Say it:} \emph{aṅgulīyakam}
  \item \emph{Say it:} \emph{aṅgulīyakam}, once more
  \item \emph{Say it:} say \emph{aṅgulīyakam}
  \item \emph{Recall:} say the Sanskrit for a ball, then the Sanskrit for a toy, then say \emph{aṅgulīyakam} again
\end{itemize}

\begin{tcolorbox}[breakable,colback=teal!4,colframe=teal!35!black,title={Before you move on}]
What does \sk{अंगुलीयकम्} mean? (\textquotedblleft{}A ring\textquotedblright{}.) Say it once more.
\end{tcolorbox}
\section[\texorpdfstring{\sk{कंकणम्} (kaṅkaṇam)}{कंकणम् (kaṅkaṇam)}]{\sk{कंकणम्} (kaṅkaṇam) — a bangle}

\label{lesson:SA-C103-kankanam}



\begin{quote}
Before the new one: say the Sanskrit for a toy, then the Sanskrit for a ring.
\end{quote}

\subsection*{You'll want to know: \sk{कंकणम्}}
\textbf{\sk{कंकणम्}} — \emph{kaṅkaṇam} — \textquotedblleft{}a bangle\textquotedblright{}.

\begin{grammarlens}[title={What you've built}]
One more word for this chapter.
\end{grammarlens}

\subsection*{Your turn}
\begin{itemize}
\raggedright
  \item \emph{Say it:} \emph{kaṅkaṇam}
  \item \emph{Say it:} \emph{kaṅkaṇam}, once more
  \item \emph{Say it:} say \emph{kaṅkaṇam}
  \item \emph{Recall:} say the Sanskrit for a toy, then the Sanskrit for a ring, then say \emph{kaṅkaṇam} again
\end{itemize}

\begin{tcolorbox}[breakable,colback=teal!4,colframe=teal!35!black,title={Before you move on}]
What does \sk{कंकणम्} mean? (\textquotedblleft{}A bangle\textquotedblright{}.) Say it once more.
\end{tcolorbox}
\section[\texorpdfstring{\sk{उत्तरीयम्} (uttarīyam)}{उत्तरीयम् (uttarīyam)}]{\sk{उत्तरीयम्} (uttarīyam) — a shawl, an upper garment}

\label{lesson:SA-C103-uttariyam}



\begin{quote}
Before the new one: say the Sanskrit for a ring, then the Sanskrit for a bangle.
\end{quote}

\subsection*{You'll want to know: \sk{उत्तरीयम्}}
\textbf{\sk{उत्तरीयम्}} — \emph{uttarīyam} — \textquotedblleft{}a shawl, an upper garment\textquotedblright{}.

\begin{grammarlens}[title={What you've built}]
One more word for this chapter.
\end{grammarlens}

\subsection*{Your turn}
\begin{itemize}
\raggedright
  \item \emph{Say it:} \emph{uttarīyam}
  \item \emph{Say it:} \emph{uttarīyam}, once more
  \item \emph{Say it:} say \emph{uttarīyam}
  \item \emph{Recall:} say the Sanskrit for a ring, then the Sanskrit for a bangle, then say \emph{uttarīyam} again
\end{itemize}

\begin{tcolorbox}[breakable,colback=teal!4,colframe=teal!35!black,title={Before you move on}]
What does \sk{उत्तरीयम्} mean? (\textquotedblleft{}A shawl, an upper garment\textquotedblright{}.) Say it once more.
\end{tcolorbox}
